\centerline{\bf \Huge Мат и драка.}

\begin{flushright}
        \scriptsize{Ну и молодежь пошла! Ты ему в морду, а он сразу в драку.\\Кенпачи Зараки}
\end{flushright}

\problem{1}
Чему равно значение выражения $\dfrac{1}{2!} + \dfrac{2}{3!} + \ldots + \dfrac{2023}{2024!}$

\problem{2}
Две равных окружности касаются изнутри третьей и касаются между собой. Соединив три центра, получим треугольник с периметром, равным 18. Найдите радиус большей окружности.


\problem{3}
Прямая, проходящая через вершину $A$ и точку $E$ на стороне $BC$ прямоугольника $ABCD$, делит прямоугольник на две части: треугольник $ABE$ и трапецию $AECD$. Известно, что $S_{ABE} : S_{AECD} = 1/6$. Найдите $BE : EC.$

\problem{4}
Дольган, Марк, Максим, Артём, Басан, Соня и Эвена начали просить у Расула Владимировича подсказки по задачам. Расулу Владимировичу это надоело и он разом выдал $2024^{2024}$ подсказок. Дети попытались поделить их поровну между собой, но остались лишние. Сколько? 

\problem{5}
На сколько нулей оканчивается число $1\cdot1! + 2\cdot2! + \ldots + 99\cdot99! + 1$

\problem{6}
Числа $x, y$ и $z$ таковы, что $\dfrac{x}{y+z} + \dfrac{y}{z+x} + \dfrac{z}{x+y} = 1$ .  Какие значения может принимать выражение  $\dfrac{x^2}{y+z} + \dfrac{y^2}{z+x} + \dfrac{z^2}{x+y}$?

\problem{7}
Среди учеников Расула Владимировича прошёл шахматный турнир. Каждый участник сыграл с каждым из остальных ровно один раз. За победу давали один ЭлМШкоин, за ничью – полЭлМШкоина, за поражение не давалось ничего. Оказалось, что среди каждых трёх участников найдётся шахматист, заработавший в партиях с двумя другими ровно 1,5 ЭлМШкоина. Какое наибольшее количество учеников могло участвовать в таком турнире?

\problem{8}
Карта Элисты представляет собой квадрат 6×6 клеток. Каждая клетка – либо территория банды, либо спорная территория. Банд всего 27, а спорных территорий 9. На спорную территорию претендуют все банды по соседству и только они (то есть клетки, соседние со спорной по стороне или вершине). Может ли быть, что на каждые две спорные территории претендует разное число банд?


\newpage

\hspace{3mm}

\problem{9}
Двузначное число в сумме с числом, записанным теми же цифрами, но в обратном порядке, даёт полный квадрат. Найти все такие числа.


\problem{10}
Сколько рациональных слагаемых содержится в разложении
а) $(\sqrt{2} + \sqrt[3]{3})^{100}$;


\problem{11}
В некоторый момент угол между часовой и минутной стрелками равен $\alpha$. Через час он опять равен $\alpha$. Найдите все возможные значения $\alpha$.


\problem{12}
В пачке 20 карточек: синие, красные и желтые. Синих в шесть раз меньше, чем желтых, и красных меньше, чем желтых. Какое наименьшее количество карточек надо вытащить не глядя, чтобы среди них обязательно оказалась красная?


\problem{13}
В треугольнике $ABC \angle C = 90^{\circ}, A0, B0, C0 –$ середины сторон $BC, CA, AB$ соответственно. На отрезках AB0 и BA0 во внешнюю сторону построены как на основаниях равносторонние треугольники с вершинами $C_1, C_2$. Найдите угол $C_0C_1C_2$.

\problem{14}
Пятиугольник $ABCDE$ описан около окружности. Углы при его вершинах $A, C$ и $E$ равны $100^{\circ}$. Найдите угол $ACE$.

\problem{15}
Дан многочлен  $x(x + 1)(x + 2)(x + 3)$.  Найти его наименьшее значение.

\problem{16}
На каждой клетке шахматной доски вначале стоит по ладье. Каждым ходом можно снять с доски ладью, которая бьет нечётное число ладей. Какое наибольшее число ладей можно снять? (Ладьи бьют друг друга, если они стоят на одной вертикали или горизонтали и между ними нет других ладей.)

\problem{17}
Точки K и N расположены соответственно на сторонах $AB$ и $AC$ треугольника $ABC$, причём  $AK = BK$  и  $AN = 2NC$. В каком отношении отрезок KN делит медиану $AM$ треугольника $ABC?$

\problem{18}
Дан треугольник $ABC$ площади $1$. На медианах $AK, BL$ и $CN$ взяты точки $P, Q$ и $R$ так, что  $AP = PK,  BQ : QL = 1 : 2,  CR : RN = 5 : 4$.  Найдите площадь треугольника $PQR$.

\problem{19}
На бесконечной ленте выписаны в ряд числа. Первой идёт единица, а каждое следующее число получается из предыдущего прибавлением к нему наименьшей ненулевой цифры его десятичной записи. Сколько знаков в десятичной записи числа, стоящего в этом ряду на $9\cdot1000^{1000}$-м месте?